\begin{table}[H]
\centering
\caption{Standardized Effect Sizes for Main Outcomes}
\label{tab:sde}
\begin{minipage}{\textwidth}
\begin{threeparttable}
\begin{adjustbox}{max width=\textwidth}
\begin{tabular}{llcccccc}
\toprule
Outcome & Specification & $\hat{\beta}$ & SE & SD($Y$) & SDE & SE(SDE) & Classification \\
\midrule
\multicolumn{8}{l}{\textit{Panel A: Pooled}} \\
Program survival & Baseline DiD & -0.2853 & 0.0102 & 0.331 & -0.862 & 0.031 & Large negative \\
Annual exit rate & CIP-level DiD & 0.0463 & 0.0113 & 0.150 & 0.309 & 0.075 & Large positive \\
Log completions & Intensive margin & -0.2236 & 0.0261 & 1.377 & -0.162 & 0.019 & Large negative \\
\midrule
\multicolumn{8}{l}{\textit{Panel B: Heterogeneous (by CIP risk)}} \\
Survival (high-risk CIPs) & Cosmetology, Business & -0.2151 & 0.0099 & 0.331 & -0.650 & 0.030 & Large negative \\
Survival (low-risk CIPs) & Healthcare & -0.2625 & 0.0129 & 0.331 & -0.793 & 0.039 & Large negative \\
\bottomrule
\end{tabular}
\end{adjustbox}
\end{threeparttable}
\par\vspace{0.5em}
{\footnotesize \textit{Notes:} \textbf{Country:} United States. \textbf{Research question:} Did federal disclosure of program-level debt-to-earnings outcomes under the Gainful Employment Rule cause for-profit educational programs to close, and did the 2019 repeal reverse these closures? \textbf{Policy mechanism:} The GE Rule required all for-profit educational programs to report graduates' debt-to-earnings ratios; programs failing thresholds faced loss of federal Title IV financial aid eligibility. The rule created a public scorecard distinguishing passing from failing programs, combining regulatory threat with information disclosure. \textbf{Outcome definition:} Program survival is a binary indicator equal to one if a program (institution $\times$ CIP $\times$ award level) reports positive completions in IPEDS in a given year. Exit rate is the annual fraction of existing programs that cease reporting. Log completions is the natural log of total program completions. \textbf{Treatment:} Binary: for-profit sector status (control = 3 in IPEDS), interacted with post-2015 indicator. \textbf{Data:} IPEDS Completions Survey (c\_a) and Institution Directory (hd), 2008--2023, at the program level (institution $\times$ CIP $\times$ award level $\times$ year). Forward cohort of programs active in 2011: 19,433 for-profit and 161,976 public programs. \textbf{Method:} Two-way fixed effects (program and year FE), with treatment defined as for-profit $\times$ post-GE (2015+). Standard errors clustered at the state level. \textbf{Sample:} Programs with positive completions in 2011, tracked through 2023. For-profit defined as IPEDS control = 3; public as control = 1. Programs in 39 CIP2 codes present in both sectors. SDE $= \hat{\beta} / \text{SD}(Y)$ where SD($Y$) is the pre-treatment standard deviation. Classification refers to magnitude, not statistical significance: Large ($|$SDE$| > 0.15$), Moderate ($0.05$--$0.15$), Small ($0.005$--$0.05$), Null ($< 0.005$).
}
\end{minipage}
\end{table}
